
\documentclass{article}
\usepackage{colortbl}
\usepackage{makecell}
\usepackage{multirow}
\usepackage{supertabular}

\begin{document}

\newcounter{utterance}

\centering \large Interaction Transcript for game `bbh', experiment `logical\_deduction\_seven\_objects', episode 189 with Qwen3.6{-}27B.
\vspace{24pt}

{ \footnotesize  \setcounter{utterance}{1}
\setlength{\tabcolsep}{0pt}
\begin{supertabular}{c@{$\;$}|p{.15\linewidth}@{}p{.15\linewidth}p{.15\linewidth}p{.15\linewidth}p{.15\linewidth}p{.15\linewidth}}
    \# & \multicolumn{2}{c}{Player} && \multicolumn{2}{c}{Game Master} \\
    \hline

    \theutterance \stepcounter{utterance}  
    & & & \multicolumn{4}{p{0.6\linewidth}}{
        \cellcolor[rgb]{0.9,0.9,0.9}{
            \makecell[{{p{\linewidth}}}]{
                \texttt{\tiny{[P1$\langle$GM]}}
                \texttt{A logical deduction task which requires deducing the order of a sequence of objects.} \\
\\ 
\texttt{Q: The following paragraphs each describe a set of three objects arranged in a fixed order. The statements are logically consistent within each paragraph. In a golf tournament, there were three golfers: Amy, Eli, and Eve. Eve finished above Amy. Eli finished below Amy.} \\
\texttt{Options:} \\
\texttt{(A) Amy finished last} \\
\texttt{(B) Eli finished last} \\
\texttt{(C) Eve finished last} \\
\texttt{A: (B)} \\
\\ 
\texttt{Q: The following paragraphs each describe a set of three objects arranged in a fixed order. The statements are logically consistent within each paragraph. On a shelf, there are three books: a white book, a green book, and an orange book. The green book is to the right of the white book. The orange book is the rightmost.} \\
\texttt{Options:} \\
\texttt{(A) The white book is the leftmost} \\
\texttt{(B) The green book is the leftmost} \\
\texttt{(C) The orange book is the leftmost} \\
\texttt{A: (A)} \\
\\ 
\texttt{Q: The following paragraphs each describe a set of three objects arranged in a fixed order. The statements are logically consistent within each paragraph. On a shelf, there are three books: a red book, a gray book, and a white book. The white book is to the left of the gray book. The red book is the second from the left.} \\
\texttt{Options:} \\
\texttt{(A) The red book is the leftmost} \\
\texttt{(B) The gray book is the leftmost} \\
\texttt{(C) The white book is the leftmost} \\
\texttt{A: (C)} \\
\\ 
\texttt{Q: The following paragraphs each describe a set of seven objects arranged in a fixed order. The statements are logically consistent within each paragraph. On a shelf, there are seven books: an orange book, a red book, a green book, a white book, a purple book, a brown book, and a yellow book. The green book is to the left of the purple book. The white book is to the right of the orange book. The green book is to the right of the brown book. The yellow book is to the right of the white book. The red book is to the right of the yellow book. The brown book is to the right of the red book.} \\
\texttt{Options:} \\
\texttt{(A) The orange book is the fourth from the left} \\
\texttt{(B) The red book is the fourth from the left} \\
\texttt{(C) The green book is the fourth from the left} \\
\texttt{(D) The white book is the fourth from the left} \\
\texttt{(E) The purple book is the fourth from the left} \\
\texttt{(F) The brown book is the fourth from the left} \\
\texttt{(G) The yellow book is the fourth from the left} \\
\texttt{A:} \\
            }
        }
    }
     \\ \\

    \theutterance \stepcounter{utterance}  
    & \multicolumn{4}{p{0.6\linewidth}}{
        \cellcolor[rgb]{0.561,0.737,0.561}{
            \makecell[{{p{\linewidth}}}]{
                \texttt{\tiny{[P1$\rangle$GM]}}
                \texttt{To determine the position of the books, let's analyze the constraints step by step. We have seven positions, numbered 1 (leftmost) to 7 (rightmost).} \\
\\ 
\texttt{The books are: Orange (O), Red (R), Green (G), White (W), Purple (P), Brown (B), Yellow (Y).} \\
\\ 
\texttt{The statements are:} \\
\texttt{1. The green book is to the left of the purple book (\$G < P\$).} \\
\texttt{2. The white book is to the right of the orange book (\$O < W\$).} \\
\texttt{3. The green book is to the right of the brown book (\$B < G\$).} \\
\texttt{4. The yellow book is to the right of the white book (\$W < Y\$).} \\
\texttt{5. The red book is to the right of the yellow book (\$Y < R\$).} \\
\texttt{6. The brown book is to the right of the red book (\$R < B\$).} \\
\\ 
\texttt{Let's chain these inequalities together.} \\
\\ 
\texttt{From statements 2, 4, and 5, we have a chain:} \\
\texttt{\$O < W < Y < R\$} \\
\\ 
\texttt{From statement 6, we add Brown to the end of that chain:} \\
\texttt{\$O < W < Y < R < B\$} \\
\\ 
\texttt{From statement 3, we know Brown is to the left of Green:} \\
\texttt{\$O < W < Y < R < B < G\$} \\
\\ 
\texttt{From statement 1, we know Green is to the} \\
            }
        }
    }
    & & \\ \\

    \theutterance \stepcounter{utterance}  
    & & & \multicolumn{2}{p{0.3\linewidth}}{
        \cellcolor[rgb]{0.9,0.9,0.9}{
            \makecell[{{p{\linewidth}}}]{
                \texttt{\tiny{[GM$|$GM]}}
                \texttt{Target: (B)} \\
            }
        }
    }
    & & \\ \\

    \theutterance \stepcounter{utterance}  
    & & & \multicolumn{2}{p{0.3\linewidth}}{
        \cellcolor[rgb]{0.9,0.9,0.9}{
            \makecell[{{p{\linewidth}}}]{
                \texttt{\tiny{[GM$|$GM]}}
                \texttt{game\_result = LOSE} \\
            }
        }
    }
    & & \\ \\

\end{supertabular}
}

\end{document}
